% OUP official authoring-template LaTeX package for Briefings in Bioinformatics.
% Generated from submission_package/01_manuscript/main_manuscript.md.
% OUP guidance: use the OUP LaTeX template with the 'numbered' option for Vancouver-style citations.
\documentclass[modern,large,numbered,nocrop]{oup-authoring-template}

\usepackage{booktabs}
\usepackage{tabularx}
\usepackage{url}
\usepackage{hyperref}
\usepackage{float}
\graphicspath{{figures/}}
\setcitestyle{numbers,square,comma}
\makeatletter
\@numbibtrue
\NAT@numberstrue
\makeatother

\journaltitle{Briefings in Bioinformatics}
\pubyear{2026}
\copyrightyear{2026}
\title[Transcriptional reach of the NK serine--SM--topology axis]{Mapping the Transcriptional Reach of the Serine--Sphingomyelin--Membrane-Topology Axis of NK-Cell Immune Evasion: A Single-Cell-Informed Heterogeneous Graph Framework from Liver to Gastric Cancer}

\author[1,2]{Guohao Lyu}
\author[1,2]{Yingchun Xia}
\author[1,2]{Huichao Liu}
\author[1,2]{Xiaolei Zhu}
\author[1,2]{Shuai Yang}
\author[3,4,*]{Ailian Zhou}
\author[1,2,*]{Lichuan Gu}
\address[1]{School of Artificial Intelligence, Anhui Agricultural University, Hefei 230036, China}
\address[2]{Anhui Province Key Laboratory of Intelligent Agricultural Technology and Equipment, Anhui Agricultural University, Hefei 230036, China}
\address[3]{Agricultural Information Institute, Chinese Academy of Agricultural Sciences, Beijing 100081, China}
\address[4]{Key Laboratory of Agricultural Blockchain Application, Ministry of Agriculture and Rural Affairs, Beijing 100081, China}
\corresp[*]{Correspondence should be addressed to Lichuan Gu. Email: glc@ahau.edu.cn}
\corresp[*]{Correspondence should be addressed to Ailian Zhou. Email: zhouailian@caas.cn}

\abstract{\subheading{Motivation} A landmark wet-lab study (Zheng et al., \textit{Nat Immunol} 2023) established that tumors evade natural killer (NK) cell cytotoxicity by dysregulating serine metabolism, which depletes sphingomyelin (SM) in NK membranes, collapses membrane protrusions, and abolishes lytic immune synapse formation. However, this mechanism was demonstrated in a small number of liver cancer patients using specialized single-cell mass spectrometry and super-resolution imaging -- methods that cannot scale to cohort-level analyses or be readily tested in other cancer types. \subheading{Results} We present GC-NKGraph-Atlas, a computational framework that asks \textbf{how much} of the serine$\to$sphingomyelin$\to$membrane-topology$\to$cytotoxicity axis is transcriptionally recoverable from public tumor transcriptomes. The framework employs: (i) a reusable \textit{mechanism-card} abstraction that encodes the wet-lab mechanism as a machine-readable recipe; (ii) a single-cell-informed, NK-aware heterogeneous graph that integrates protein--protein interactions, ligand--receptor pairs, transcription factor targets, and a mechanism-grounded \texttt{metabolic\_crosstalk} edge type; and (iii) a graph neural network that learns gene embeddings from this multi-relational graph to predict NK immune states. A two-arm design (Arm A: liver/HCC positive control; Arm B: gastric cancer extension) yields a \textbf{scoping result}: the \textit{effector arm} of the axis reproduces robustly and independently, with protrusion-machinery coupled to cytotoxicity output in TCGA-LIHC (r=0.55, p=5$\times$$10^{-35}$) and in 8,310 single NK cells (r=0.32, p<$10^{-193}$), and intratumoral NK cells show the expected loss of cytotoxic output ($\Delta$=$-$0.14, p=6$\times$$10^{-52}$). The \textit{upstream metabolic arm} is only weakly and partially recoverable: SM-balance$\to$protrusion coupling is undetectable in bulk but becomes statistically detectable after cell-type resolution (single-cell NK r=0.030, p=6$\times$$10^{-3}$), whereas topology-machinery transcription does \textbf{not} track the physical topology phenotype (intratumoral NK show \textit{higher}, not lower, protrusion-machinery transcript levels). This delineates which layers of a physical/metabolic immune-evasion mechanism are transcriptionally accessible, and the recovered effector coupling replicates in two independent gastric microarray cohorts (GSE62254 r=0.42; GSE84437 r=0.62; both p$\ll$$10^{-13}$). The framework then prioritizes \textbf{37 putative tumor-intrinsic candidate targets} (led by the druggable serine/sphingomyelin enzymes PHGDH, SGMS2, PSAT1, PSPH, SMPD3/1), each with a recommended wet-lab validation assay, kept strictly separate from the NK-side axis readout panel; we note that their malignant-cell fold-changes are near-zero and $\sim$46\% remain annotated to NK-side modules, so this list is a mechanism-intersecting shortlist for follow-up rather than a set of tumor-exclusive targets. \subheading{Availability} All code, configuration, and synthetic test data are available at \url{https://github.com/nblvguohao/GC-NKGraph-Atlas}. The mechanism-card template enables application to additional published mechanisms without modifying the pipeline core.}
\keywords{NK cell; immune evasion; heterogeneous graph neural network; single-cell transcriptomics; sphingomyelin; gastric cancer}

\begin{document}
\maketitle

\section*{Key Points}
\begin{itemize}
\item A reusable \textbf{mechanism-card} formalism converts published wet-lab immune-evasion mechanisms into scalable transcriptome-based target-discovery runs.
\item A \textbf{two-arm design} (liver positive control + gastric extension) produces a \textbf{scoping result} rather than a blanket claim: it separates the parts of a physical immune-evasion mechanism that are strongly transcriptionally recoverable (the protrusion$\to$cytotoxicity effector arm) from those that are weak, partial, or not recoverable (the metabolic arm and the physical topology phenotype itself).
\item \textbf{Cell-type resolution is decisive:} the SM-balance$\to$protrusion coupling is invisible in bulk but weakly detectable in single NK cells, supporting single-cell attribution rather than bulk-only inference for this class of mechanism.
\item The heterogeneous graph introduces a \textbf{mechanism-grounded \texttt{metabolic\_crosstalk} edge} (tumor serine program $\to$ NK topology state) justified by biology, not generic priors.
\item The framework outputs a \textbf{de-circularized, putative tumor-intrinsic candidate target list} (37 genes, led by druggable serine/SM enzymes) kept strictly separate from the NK-side axis readout, each with a recommended wet-lab validation assay -- with the caveat that the candidates carry near-zero malignant-cell fold-changes and $\sim$46\% retain NK-side module annotation, so the list is a mechanism-intersecting shortlist, not a set of tumor-exclusive targets.
\end{itemize}

\section{Introduction}

\subsection{Biological context}

Natural killer (NK) cells are innate lymphoid cells critical for anti-tumor immunity. Unlike T cells, NK cells kill without prior antigen sensitization, making them attractive effectors for cancer immunotherapy [1,2,25--27]. However, tumors deploy multiple mechanisms to evade NK-mediated killing: metabolic competition [3,4,45], inhibitory ligand upregulation [5], immunosuppressive cytokine secretion [6], and physical exclusion from the tumor nest [7].

A particularly elegant evasion mechanism was recently elucidated by Zheng et al. [8]: tumors dysregulate serine metabolism [28,29], reducing the availability of serine for NK sphingolipid synthesis. The resulting depletion of sphingomyelin [30,31] in NK membranes collapses membrane protrusions and microvilli, preventing the formation of lytic immune synapses and abolishing cytotoxicity. Critically, this phenotype can be rescued by inhibiting SM catabolism (via ASM/SMPD1 or NSMASE/SMPD2-4), and the rescue is synergistic with Tim3 (HAVCR2) checkpoint blockade. The follow-up framing [9] positions sphingomyelin as a metabolic immune checkpoint and reports that the mechanism extends beyond liver cancer to lung, colon, and ovarian cancers.

\textbf{However, this mechanism was proven with techniques that do not scale.} The anchor paper used super-resolution SEM imaging and single-immunocyte mass spectrometry on a small patient cohort -- gold-standard evidence for physical membrane topology but limited to specialized laboratories and small sample sizes. No transcriptomic data were deposited. The transcriptional footprint of this mechanism -- i.e., whether the molecular machinery and capacity for the serine$\to$SM$\to$protrusion$\to$cytotoxicity axis can be detected from widely available transcriptomic data -- remains unexplored.

\subsection{The gap}

Three gaps motivate this work:

\begin{enumerate}
\item \textbf{Scalability.} The mechanism was demonstrated in tens of patients; testing it across hundreds or thousands of tumor samples requires a transcriptome-based proxy.
\item \textbf{Transferability.} Whether the same axis operates in gastric cancer -- a digestive-tract cancer not on the published extension list (lung/colon/ovarian) [9] -- is unknown.
\item \textbf{Reusability.} Each published mechanistic discovery currently requires a bespoke computational follow-up; a generalizable engine that converts mechanism papers into target-discovery runs does not exist.
\end{enumerate}

\subsection{Our approach}

We introduce \textbf{GC-NKGraph-Atlas}, a computational framework designed to address all three gaps. The framework is organized around a reusable \textbf{mechanism-card} abstraction: a machine-readable YAML specification that encodes one published wet-lab mechanism -- its molecular chain, gene modules, expected directions, physical ground-truth targets, validation assays, and therapeutic hooks -- as a recipe that the pipeline consumes.

The specific mechanism card driving this study (\texttt{zheng\_nk\_sm\_topology}) operationalizes the Zheng 2023 serine--SM--topology axis. The framework: (a) defines cell-type-attributed transcriptional proxies for each step of the mechanistic chain, computed from single-cell RNA-seq; (b) constructs a tumor--NK heterogeneous graph in which a \texttt{metabolic\_crosstalk} edge type is justified by the specific biology rather than generic interaction priors; (c) learns gene embeddings from this graph to predict NK immune states; and (d) ranks putative tumor-intrinsic targets by multi-evidence scoring.

The study employs a two-arm design:

\begin{itemize}
\item \textbf{Arm A (Positive Control, liver/HCC):} test how much of the published axis is visible in the system where it was proven, using independent public transcriptomes (TCGA-LIHC + public HCC scRNA). This is the credibility anchor.
\item \textbf{Arm B (Novel Extension, gastric cancer):} test the same axis in gastric cancer -- a natural, not-yet-claimed digestive-tract extension -- and prioritize gastric-specific candidate targets.
\end{itemize}

\subsection{Related work}

Our framework draws on four strands of computational biology, and its novelty lies in how it combines them around a specific published mechanism rather than in any single component.

\textbf{Immune-state and cell-type inference from transcriptomes.} A large tool family scores immune infiltration and cell state from bulk and single-cell data -- CIBERSORTx [17] and quanTIseq [18] for deconvolution, and phenotype-to-genotype methods such as Scissor [19], which links single-cell phenotypes to bulk clinical variables. These estimate \textit{how much} of a cell type is present or \textit{which} phenotype dominates; none operationalizes a defined mechanistic chain or enforces the claim boundary between a transcriptional proxy and a physical phenotype, which is central to our design. Single-cell atlases increasingly link a functional cell state to long-term clinical outcome -- for example, a type-2 program in CAR-T infusion products associated with multi-year leukaemia remission [46]; our framework instead attributes states to a defined evasion mechanism, and anchoring the NK-state readout to clinical outcome remains future work (\S{}4.3).

\textbf{Cell--cell communication and metabolic-flux inference.} CellChat [20], CellPhoneDB [16] and NicheNet [21] infer ligand--receptor signaling between cell types, and scFEA [22] estimates metabolic flux from single-cell expression. We borrow the ligand--receptor edge concept for the heterogeneous graph but add a mechanism-specific \texttt{metabolic\_crosstalk} edge grounded in one published serine$\to$sphingomyelin relationship, rather than a generic communication prior.

\textbf{Trajectory and dysfunction-state modeling.} CytoTRACE [23] and related pseudotime methods order cells along differentiation or dysfunction gradients; we use an analogous NK dysfunction axis to define reversible-state proxy labels, but anchor the states to the specific effector/checkpoint genes of the target mechanism rather than to a generic exhaustion signature. Most characterized regulators of NK dysfunction are surface checkpoints; intracellular and transcription-factor-level regulators (e.g. the cAMP-responsive factor CREM downstream of CAR and IL-15 signalling [47]) are comparatively understudied, and our candidate prioritization currently inherits this surface/metabolic-gene bias (\S{}4.3).

\textbf{Graph learning in cancer genomics.} Graph neural networks [36--38] have been applied to multi-omics integration (e.g. MOGONET [24]) and to molecular-interaction networks for outcome prediction, and heterogeneous graph transformers [10] provide type-specific message passing over multi-relational graphs. Existing applications typically treat the graph as a generic prior (a PPI or co-expression network). Our contribution is a graph whose \textit{every edge type carries an explicit biological justification}, including one edge derived directly from the mechanism under study, so that the heterogeneous architecture has a principled reason to exist for this problem.

To our knowledge, no prior tool converts a single published wet-lab immune-evasion mechanism into a machine-readable card that drives cell-type-attributed proxy construction, mechanism-grounded graph learning, and de-circularized target prioritization, while explicitly measuring -- rather than assuming -- how much of the mechanism the transcriptome can reach.

\section{Methods}

\subsection{Study design overview}

The framework proceeds through 14 phases, grouped into five stages:

\begin{table*}[htbp]
\caption{GC-NKGraph-Atlas pipeline stages.}
\small
\begin{tabularx}{\textwidth}{@{}lX@{}}
\toprule
Stage & Content \\
\midrule
I -- DATA (Phases 1--2) & Download + preprocess TCGA-STAD, TCGA-LIHC, GEO gastric cohorts \\
II -- scRNA (Phases 3--7) & scRNA integration, NK atlas annotation, state scoring, trajectory \\
III -- GRAPH (Phase 8) & Heterogeneous gene graph with mechanism-grounded edges \\
IV -- MODEL (Phases 9--10) & Baseline comparison + GNN-based NK-state classifier \\
V -- TARGETS (Phases 11--14R) & SST-axis scoring, candidate prioritization, assay recommendation \\
\bottomrule
\end{tabularx}
\end{table*}

A master pipeline launcher (\texttt{src/pipeline.py}) orchestrates execution with checkpoint-based skipping and supports synthetic data mode for testing.

\begin{figure*}[htbp]
\centering
\includegraphics[width=\textwidth]{figures/fig0_workflow.pdf}
\caption{Workflow overview of the GC-NKGraph-Atlas framework, from data acquisition through target prioritization, showing the mechanism-card abstraction and the two-arm study design (Arm A: liver positive control; Arm B: gastric cancer extension). \textit{Alt text: Workflow overview of the GC-NKGraph-Atlas framework showing the computational pipeline stages, the mechanism-card abstraction, and the two-arm study design.}}\label{fig:workflow}
\end{figure*}

\subsection{The mechanism-card abstraction}

A \textit{mechanism card} is a YAML document with the following required sections (see \texttt{configs/mechanism\_cards/mechanism\_card.template.yaml} for the full schema):

\begin{table*}[htbp]
\caption{Mechanism-card schema.}
\small
\begin{tabularx}{\textwidth}{@{}lX@{}}
\toprule
Section & Content \\
\midrule
\texttt{origin} & Paper DOI, lab, notes on data availability \\
\texttt{biology} & Phenotype, cell type affected, ordered mechanistic chain \\
\texttt{transcriptional\_proxy} & Gene modules per step, expected direction, attribution requirements, claim boundaries \\
\texttt{physical\_ground\_truth} & Gated physical measurement targets (SEM, MS), mock policy \\
\texttt{graph\_integration} & New node/edge types, construction rules \\
\texttt{validation} & Positive control cohort, pre-registered hypotheses, recovery definition \\
\texttt{therapeutic\_hook} & Intervention logic, combination rationale, patient stratification readout \\
\bottomrule
\end{tabularx}
\end{table*}

\textbf{Design principle.} The card strictly separates what is \textit{computable} from transcriptomes (transcriptional proxy, status ACTIVE) from what requires \textit{physical measurement} (membrane protrusion density, SM metabolite content -- status GATED). The loader raises a clear error if gated targets are absent; mock data are permitted only for code-path testing with \texttt{*\_MOCK\_*} file patterns and are excluded from all scientific claims.

\subsection{SST-axis transcriptional proxy (Phase 14R)}

The serine--sphingomyelin--topology axis is operationalized as seven gene modules derived from the anchor paper [8] and its follow-up [9]:

\begin{table*}[htbp]
\caption{SST-axis transcriptional proxy modules.}
\footnotesize
\begin{tabularx}{\textwidth}{@{}lXXX@{}}
\toprule
Module & Cell type & Genes (n) & Expected direction \\
\midrule
\texttt{tumor\_serine\_capacity} & malignant & PHGDH, PSAT1, PSPH, SHMT1/2, MTHFD1/2, MTHFD1L, SLC1A4/5 (10) & Calibrated on liver control \\
\texttt{nk\_sm\_synthesis} & NK & SGMS1, SGMS2 (2) & $\uparrow$ = more topology-permissive \\
\texttt{nk\_sm\_catabolism} & NK & SMPD1--4 (4) & $\uparrow$ = less topology-permissive \\
\texttt{nk\_denovo\_sphingolipid} & NK & SPTLC1--3, SPTSSA, CERS2/4/5/6, DEGS1 (9) & Context-dependent \\
\texttt{nk\_protrusion\_machinery} & NK & ERM family, Arp2/3 complex, WASP/WAVE, Rho GTPases, formins, BAR domain (25) & $\uparrow$ = more topology-permissive \\
\texttt{nk\_synapse\_cytotoxicity\_outcome} & NK & NKG7, GNLY, GZMB, PRF1, IFNG, LCP2, LAT, VAV1, TLN1, ITGAL, ITGB2 (11) & Axis-positive correlate \\
\texttt{checkpoint\_link} & NK & HAVCR2 (1) & $\uparrow$ = less topology-permissive \\
\bottomrule
\end{tabularx}
\end{table*}

Per-cell module scores are computed as mean z-scores of constituent genes. Derived scores include:

\begin{itemize}
\item \textbf{nk\_sm\_balance} = mean\_zscore(sm\_synthesis) $-$ mean\_zscore(sm\_catabolism)
\item \textbf{nk\_topology\_permissive} = composite(sm\_balance, protrusion\_machinery)
\item \textbf{sst\_axis\_score} = integrated(tumor\_serine\_capacity [calibrated sign], nk\_topology\_permissive, cytotoxicity\_outcome)
\end{itemize}

\textbf{Cell-type attribution.} All tumor-vs-NK axis claims require single-cell resolution. For bulk data, deconvolution (CIBERSORTx or quanTIseq) is applied as a fallback, and results lacking cell-type resolution are flagged \texttt{MIXED\_UNRESOLVED}.

\textbf{Honesty rule.} The permitted language is strictly "transcriptional program permissive-of / associated-with the topology phenotype." Claims of physical membrane topology prediction or SM metabolite content prediction from transcriptomes are prohibited.

\subsection{scRNA-seq analysis}

\textbf{Data and loading.} Single-cell RNA-seq was obtained from GSE246662, comprising nine samples across three tissues -- healthy liver (HL1--3), gastric cancer (GC1--3), and gastric-cancer liver metastasis (LM1--3) -- totaling 166,829 cells. Per-sample count matrices were loaded with orientation auto-detection (genes$\times$cells vs. cells$\times$genes), tagged with tissue and tumor/normal condition, and concatenated on the intersection of genes (inner join) to avoid cross-platform sparsity artifacts (\texttt{src/scrna\_analysis/run\_scrna\_v2.py}).

\textbf{Quality control.} Cells were filtered on standard criteria (\texttt{src/scrna\_analysis/qc\_filter.py}): $\geq$ 200 and $\leq$ 6,000 detected genes per cell (the upper bound as a multiplet proxy) and $\leq$ 20\% mitochondrial counts; genes detected in fewer than 3 cells were removed. Mitochondrial fraction was computed with a NaN-robust fallback to handle heterogeneous cross-platform inputs. Doublet removal (Scrublet) is applied optionally per sample. QC thresholds are logged and reported in \texttt{results/tables/scrna\_qc\_summary.tsv}.

\textbf{Normalization, integration, and clustering.} Counts were library-size normalized to $10^{4}$ and log1p-transformed; 3,000 highly variable genes were selected (Seurat v3 flavor [33], with a variance-based fallback; Seurat v4 [32] provides the broader multimodal framework). Batch effects across the nine samples were corrected with scVI [14] (\texttt{sample\_id} as batch key, 30 latent dimensions, 2 layers, up to 200 epochs with early stopping). Neighborhood graphs, UMAP, and Leiden clustering (resolution 1.0) were computed on the scVI latent space using SCANPY [13]. The analysis pipeline follows current best-practice recommendations for single-cell RNA-seq [34].

\textbf{Cell-type annotation.} Lineages were assigned from canonical marker-set mean expression: NK (NCAM1, KLRD1, NKG7, GNLY, KLRF1, EOMES, NCR1, FCGR3A), T (CD3D/E/G, CD4, CD8A), monocyte (CD14, CD68, CSF1R), and B (MS4A1, CD79A, CD19). NK cells were separated from T cells by requiring an NK score above threshold with a low T score, yielding 8,310 NK cells used for the axis analyses. We note this marker-threshold labeling as a limitation (\S{}4.3) relative to reference-based mapping (e.g. scANVI/scArches [35]), and report NK counts per sample (\texttt{results/tables/gc\_scrna\_dataset\_summary.tsv}) for transparency.

\textbf{NK immune state classification.} Four states are defined from scRNA signatures:

\begin{itemize}
\item \textbf{NK-hot-cytotoxic:} high GZMB/PRF1/IFNG, low HAVCR2/TIGIT
\item \textbf{NK-hot-dysfunctional:} high HAVCR2/TIGIT/CD96, moderate cytotoxicity genes
\item \textbf{NK-cold/excluded:} low NK signature, low cytotoxicity
\item \textbf{NK-intermediate:} transitional state
\end{itemize}

These states are projected onto bulk samples via scRNA-anchored scoring (mean z-score of state-specific gene sets [43,44], calibrated on the scRNA-defined states).

\subsection{Heterogeneous graph construction}

The heterogeneous gene graph integrates six edge types across multiple node types:

\begin{table*}[htbp]
\caption{Heterogeneous graph edge types.}
\footnotesize
\begin{tabularx}{\textwidth}{@{}lXXX@{}}
\toprule
Edge type & Source $\to$ Target & Source database & Weight \\
\midrule
\texttt{ppi} & gene $\leftrightarrow$ gene & STRING v12 (score $\geq$ 700) & normalized score / 1000 \\
\texttt{ligand\_receptor} & gene $\to$ gene & CellChatDB & 0.9 \\
\texttt{tf\_target} & TF gene $\to$ target gene & ChEA 2022 & 0.8 \\
\texttt{metabolic\_crosstalk} & tumor\_serine gene $\to$ NK topology gene & Zheng 2023 & 0.5 \\
\texttt{sm\_topology\_axis} & NK axis gene $\leftrightarrow$ NK axis gene & Zheng 2023 & 0.3 \\
\texttt{dysfunction\_correlation} & gene $\to$ NK state node & Bulk correlation & abs(corr) \\
\bottomrule
\end{tabularx}
\end{table*}

\textbf{Key design choice -- the \texttt{metabolic\_crosstalk} edge.} This edge type connects tumor-side serine metabolism genes (PHGDH, PSAT1, etc.) to NK-side SM/topology genes (SGMS1, SMPD1, EZR, etc.). Unlike generic co-expression edges, this edge is specifically justified by the Zheng 2023 mechanism. The sign of the crosstalk is \textit{calibrated} on the liver positive-control cohort, not hard-coded -- a critical distinction that prevents circular reasoning.

\subsection{Graph neural network model}

The model employs a two-stage architecture:

\textbf{Stage 1 -- Gene Graph Encoder.} Gene embeddings are learned from the heterogeneous graph via spectral decomposition of the combined normalized adjacency matrix (PPI + LR + TF + SST edges). Each gene is represented as a d-dimensional vector that encodes its multi-relational neighborhood. When \texttt{torch\_geometric} [15] is available, a heterogeneous graph transformer (HGT) [10] is used instead for learnable message-passing with type-specific projections.

\textbf{Stage 2 -- NK State Classifier.} For each bulk tumor sample with expression vector $x \in \mathbb{R}^{|G|}$, the input to the classifier is the concatenation $[x, x \cdot E]$ where $E \in \mathbb{R}^{|G| \times d}$ is the gene embedding matrix. This provides both the raw expression signal and the graph-informed projection. A 2-3 layer MLP with batch normalization and dropout predicts NK immune state (binary: cytotoxic vs rest).

Training uses 5-fold stratified cross-validation with early stopping (patience=30 epochs). The optimizer is Adam (lr=1e-3, weight\_decay=1e-5) with CrossEntropy loss.

\subsection{Baseline methods}

Six baselines are evaluated on the same data splits: XGBoost [11], LightGBM [12], Random Forest, ElasticNet logistic regression, RBF SVM, and a 2-layer MLP. All baselines use the full gene expression matrix without graph structure. The comparison isolates the contribution of graph-informed features.

\subsection{Candidate target prioritization}

Genes are ranked by a composite score integrating five evidence dimensions:

\begin{table*}[htbp]
\caption{Candidate-prioritization evidence dimensions.}
\small
\begin{tabularx}{\textwidth}{@{}lXX@{}}
\toprule
Dimension & Weight & Source \\
\midrule
Tumor cell specificity (log2 FC malignant vs non-malignant) & 0.30 & scRNA \\
NK dysfunction correlation (abs) & 0.20 & scRNA NK subset \\
SST-axis membership & 0.30 & Mechanism card \\
Axis core membership (protrusion$\to$cytotoxicity) & 0.10 & Mechanism card \\
Gold-standard literature support & 0.10 & Curated list \\
\bottomrule
\end{tabularx}
\end{table*}

Each top-ranked candidate receives a recommended wet-lab validation assay (e.g., NK-tumor co-culture cytotoxicity assay, SEM membrane-protrusion imaging, sphingomyelinase-inhibitor rescue $\pm$ Tim3 blockade).

\subsection{Pre-registered hypotheses (Arm A -- liver positive control)}

The following hypotheses are registered before execution:

\begin{table*}[htbp]
\caption{Pre-registered hypotheses.}
\small
\begin{tabularx}{\textwidth}{@{}lXX@{}}
\toprule
ID & Hypothesis & Expected direction \\
\midrule
H1 & tumor\_serine\_capacity $\perp$ nk\_sm\_balance & Direction calibrated, reported \\
H2 & nk\_sm\_balance (+) nk\_protrusion\_machinery & Positive \\
H3 & nk\_protrusion\_machinery (+) cytotoxicity anchors & Positive \\
H4 & nk\_topology\_permissive ($-$) HAVCR2 / dysfunction & Negative \\
H5 & intratumoral NK < peritumoral/peripheral NK in sm\_balance \& protrusion machinery (scRNA) & Negative (tumor < normal) \\
\bottomrule
\end{tabularx}
\end{table*}

\textbf{Recovery definition.} The pre-registered full-recovery criterion required H2--H5 to pass in the expected direction in liver cancer. Because this criterion was not met, Arm A is interpreted as a partial-recovery scoping test: it identifies the effector layer that is recoverable from transcriptomes, the metabolic coupling that becomes detectable only after cell-type resolution, and the physical topology layer that remains outside transcriptional reach.

\section{Results}

\subsection{Dataset summary}

Four independent bulk cohorts and one multi-tissue scRNA dataset were processed (Table~\ref{tab:datasets}). TCGA-STAD represents the gastric adenocarcinoma molecular landscape [39], while GSE62254 provides the ACRG molecular subtype annotation [40]. Bulk NK immune-state labels were assigned by scRNA-anchored scoring; the two external microarray cohorts required probe-to-symbol remapping before NK markers could be scored (see Methods; external-validation scoring is reported where gene coverage permits).

\begin{table*}[htbp]
\caption{Dataset characteristics.}\label{tab:datasets}
\footnotesize
\begin{tabularx}{\textwidth}{@{}lXXXX@{}}
\toprule
Dataset & Cancer & Type & Samples / cells & Role \\
\midrule
TCGA-LIHC & Liver & Bulk RNA-seq & 423 & Positive control (Arm A) \\
TCGA-STAD & Gastric & Bulk RNA-seq & 450 & Train primary (Arm B) \\
GSE62254 (ACRG) & Gastric & Bulk microarray & 300 & External validation \\
GSE84437 & Gastric & Bulk microarray & 483 & External validation \\
GSE246662 & Liver/gastric/metastasis & scRNA-seq & 166,829 cells (8,310 NK) & NK atlas + axis test \\
\bottomrule
\end{tabularx}
\end{table*}

\subsection{Arm A -- the axis is \textit{partially} recovered: effector arm yes, topology arm no}

We tested the pre-registered hypotheses (H1--H5) at two resolutions -- bulk TCGA-LIHC and single NK cells -- and report the outcome per hypothesis (Table~\ref{tab:hyp-outcomes}; \texttt{results/tables/sst\_axis\_positive\_control\_recovery.tsv}). The picture is consistent and interpretable rather than a uniform pass:

\begin{itemize}
\item \textbf{The metabolic-serine capacity is transcriptionally uncoupled from NK SM balance (H1).} Tumor-side serine-pathway transcript scores show no correlation with NK-side SM-balance scores at either resolution (bulk r=$-$0.016, $P$=0.74; scNK r=+0.012, $P$=0.27). This is consistent with the mechanism's core premise that serine$\to$SM crosstalk operates at the metabolite level -- enzyme abundance (transcription) is not flux -- and provides an empirical null baseline for the axis hierarchy.
\item \textbf{The effector arm reproduces robustly (H3).} Protrusion-machinery transcript score couples to the cytotoxicity-output score both in bulk TCGA-LIHC (r=0.551, $P$=5.5$\times$$10^{-35}$) and across 8,310 single NK cells (r=0.318, $P$=4.6$\times$$10^{-194}$), in every tissue examined. This is an independent replication of the axis's functional endpoint from data the anchor lab never deposited. To rule out the concern that this coupling is merely co-activation of two facets of the NK synapse program, we computed a generic NK-activation signature from 16 classic markers (CD69, TNF, XCL1/2, CCL3/4/5, CSF2, IL2RA, ICOS, TNFSF10, FASLG, CD38, HLA-DRA/B1, MKI67) not present in either module, then computed the partial correlation controlling for activation. The protrusion$\to$cytotoxicity coupling remains robust (partial r=0.286, $P$=1.6$\times$$10^{-156}$; r$^{2}$=0.082 vs raw r$^{2}$=0.122), with only 33\% of the shared variance attributable to co-activation. The residual is a real, independent effect -- the effector-arm recovery is not an artifact of module overlap. That tumor-infiltrating NK co-express activation and dysfunction programs is well documented [47], which is precisely why we control for a generic activation signature here rather than reporting the raw coupling alone.
\item \textbf{The upstream metabolic coupling is statistically detectable but of negligible magnitude (H2).} The SM-balance$\to$protrusion coupling is undetectable in bulk (r=$-$0.017, $P$=0.72) and becomes statistically significant after NK-cell isolation (r=+0.030, $P$=6$\times$$10^{-3}$, $P_{\mathrm{FDR}}$=9$\times$$10^{-3}$; r$^{2}$=0.0009, 95\% CI [0.009, 0.052]), but the shared variance is under 0.1\% -- the effect is statistically distinguishable from zero but biologically negligible. This does not rescue the metabolic-arm claim; it shows that cell-type resolution reveals a tiny per-cell signal that bulk averaging entirely submerges.
\item \textbf{The physical topology phenotype is \textit{not} captured by machinery transcription (H4, H5).} The topology-permissive$\to$dysfunction relationship carries the wrong sign at both resolutions (bulk r=+0.311; single-cell vs HAVCR2 r=+0.050), and intratumoral NK cells show \textit{higher} protrusion-machinery transcription than normal-tissue NK ($\Delta$=+0.142, $P$=3.0$\times$$10^{-91}$) -- the opposite of the physical collapse -- even though their cytotoxic output is correctly reduced ($\Delta$=$-$0.141, $P$=5.9$\times$$10^{-52}$). Transcription of the machinery genes therefore does not proxy the membrane-topology state itself. The reduced cytotoxic output of intratumoral NK (H5-cytotoxicity) is independently corroborated by a single-cell HCC study reporting that intratumoral, relative to peritumoral, NK cells upregulate inhibitory-checkpoint and exhaustion programs and downregulate cytotoxicity pathways [48] -- external support, from cohorts independent of ours, for the effector-layer direction; the opposing protrusion-transcription result remains our own observation.
\end{itemize}

\begin{table*}[htbp]
\caption{Pre-registered hypothesis outcomes (multi-resolution). $P_{\mathrm{FDR}}$: Benjamini--Hochberg false discovery rate correction across 10 tests.}\label{tab:hyp-outcomes}
\scriptsize
\resizebox{\textwidth}{!}{\begin{tabular}{@{}p{0.7cm}p{2.2cm}p{1.3cm}p{0.8cm}p{0.9cm}p{1.0cm}p{0.8cm}p{1.8cm}@{}}
\toprule
Hyp & Test & Resolution & r / $\Delta$ & $P$ & $P_{\mathrm{FDR}}$ & Expected & Outcome \\
\midrule
H1 & serine\_cap \textasciitilde{} sm\_balance & bulk & $-$0.016 & 0.74 & 0.74 & calibrated & reported (null) \\
H1 & serine\_cap \textasciitilde{} sm\_balance & single-cell NK & +0.012 & 0.27 & 0.34 & calibrated & reported (null) \\
H2 & sm\_balance \textasciitilde{} protrusion & bulk & $-$0.017 & 0.72 & 0.74 & + & not recovered \\
H2 & sm\_balance \textasciitilde{} protrusion & single-cell NK & +0.030 & 6$\times$$10^{-3}$ & 9$\times$$10^{-3}$ & + & stat.\ detectable, negligible \\
H3 & protrusion \textasciitilde{} cytotoxicity & bulk & +0.551 & 5$\times$$10^{-35}$ & 1$\times$$10^{-34}$ & + & \textbf{recovered} \\
H3 & protrusion \textasciitilde{} cytotoxicity & single-cell NK & +0.318 & 5$\times$$10^{-194}$ & 5$\times$$10^{-193}$ & + & \textbf{recovered} \\
H4 & topology \textasciitilde{} dysfunction & bulk & +0.311 & 7$\times$$10^{-11}$ & 4$\times$$10^{-5}$ & $-$ & not recovered \\
H4 & topology \textasciitilde{} HAVCR2 & single-cell NK & +0.050 & 4$\times$$10^{-6}$ & 8$\times$$10^{-6}$ & $-$ & not recovered \\
H5 & intratumoral<normal: cytotoxicity & single-cell NK & $-$0.141 & 6$\times$$10^{-52}$ & 9$\times$$10^{-46}$ & tumor< & \textbf{recovered} \\
H5 & intratumoral<normal: protrusion & single-cell NK & +0.142 & 3$\times$$10^{-91}$ & 1$\times$$10^{-84}$ & tumor< & not recovered \\
\bottomrule
\end{tabular}}
\end{table*}

\textbf{Recovery verdict (revised).} H1 is null at both resolutions, establishing that tumor serine-pathway transcription and NK SM balance are empirically uncoupled at the transcript level -- consistent with the serine$\to$SM link operating post-transcriptionally. Arm A recovers the \textit{functional/effector layer} of the axis (H3 across resolutions; H5-cytotoxicity) and detects a \textit{negligible-magnitude metabolic coupling} (H2 in NK cells, r$^{2}$<0.001), but does \textbf{not} recover the \textit{physical topology phenotype} from machinery transcription (H4; H5-protrusion). This is the paper's central finding -- a map of the mechanism's transcriptional reach, from null (H1) to strong (H3) to absent (H4/H5-protrusion).

\begin{figure*}[htbp]
\centering
\includegraphics[width=\textwidth]{figures/fig1_armA_positive_control.pdf}
\caption{(A) Per-sample SST module scores in TCGA-LIHC. (B) H2/H3 at bulk vs single-cell resolution (the resolution-rescue effect). (C) Intratumoral vs normal NK for cytotoxicity-output and protrusion-machinery. (D) Table~\ref{tab:hyp-outcomes} as a forest plot. \textit{Alt text: Multi-panel summary of liver positive-control recovery, showing SST module scores, bulk and single-cell hypothesis tests, intratumoral versus normal NK-cell comparisons, and the hypothesis-outcome forest plot.}}\label{fig:armA}
\end{figure*}

\subsection{Arm B -- Gastric cancer extension}

The effector coupling that defines the recovered arm holds in gastric-cancer NK cells (protrusion\textasciitilde{}cytotoxicity r=0.346, p=6.6$\times$$10^{-30}$, n=1,017 NK), i.e. the transcriptionally-recoverable layer of the axis extends to gastric cancer. Cross-tissue module means (\texttt{sst\_axis\_scrna\_by\_tissue.tsv}) show gastric-cancer NK with the lowest cytotoxicity-output score ($-$0.194) of the three tissues, consistent with an evasion phenotype. Gastric bulk (TCGA-STAD) NK states distribute across cold/excluded (154), hot-cytotoxic (153), intermediate (123), and hot-dysfunctional (20).

\textbf{External validation in two independent gastric cohorts.} The external microarray cohorts initially failed NK scoring because probe IDs were not mapped to gene symbols; after platform-annotation remapping (GSE62254$\to$GPL570, 54,675 probes$\to$22,880 genes, NK markers 6/7; GSE84437$\to$GPL6947, 49,576$\to$25,159 genes, NK markers 7/7; \texttt{run\_geo\_external\_validation.py}), the recovered effector coupling replicates in \textbf{both} cohorts: protrusion-machinery$\to$cytotoxicity-output correlates at r=0.42 (p=1.4$\times$$10^{-14}$) in GSE62254 (n=300) and r=0.62 (p=3.3$\times$$10^{-53}$) in GSE84437 (n=483) -- matching the direction and comparable in strength to the liver control (Table~\ref{tab:extval}). Notably, the SM-balance$\to$protrusion coupling that was undetectable in bulk \textit{liver} is weakly but significantly positive in both bulk \textit{gastric} cohorts (r=0.18, p=1.3$\times$$10^{-3}$; r=0.11, p=0.02), consistent with the effect being real but small and easier to detect where NK signal is stronger. The effector layer of the axis is therefore independently reproduced in three gastric datasets (one scRNA, two bulk microarray).

\begin{table*}[htbp]
\caption{External validation of the axis (independent gastric cohorts).}\label{tab:extval}
\footnotesize
\begin{tabularx}{\textwidth}{@{}lXXXX@{}}
\toprule
Cohort & n & protrusion\textasciitilde{}cytotoxicity r (p) & sm\_balance\textasciitilde{}protrusion r (p) & NK markers \\
\midrule
GSE62254 (GPL570) & 300 & 0.42 (1.4$\times$$10^{-14}$) & 0.18 (1.3$\times$$10^{-3}$) & 6/7 \\
GSE84437 (GPL6947) & 483 & 0.62 (3.3$\times$$10^{-53}$) & 0.11 (0.02) & 7/7 \\
\bottomrule
\end{tabularx}
\end{table*}

\begin{figure*}[htbp]
\centering
\includegraphics[width=\textwidth]{figures/fig2_armB_extension.pdf}
\caption{(A) SST-axis scores in TCGA-STAD; (B) cross-tissue module comparison (liver, gastric, metastasis) with gastric-cancer NK showing the lowest cytotoxicity-output score; (C) effector coupling replication across four cohorts (TCGA-LIHC, TCGA-STAD, GSE62254, GSE84437). \textit{Alt text: Multi-panel gastric-cancer extension figure. (A) SST-axis scores. (B) Tissue-level module comparisons. (C) Effector-coupling replication across four independent cohorts.}}\label{fig:armB}
\end{figure*}

\subsection{NK immune state classification}

On TCGA-STAD (5-fold stratified CV, binary NK-hot-cytotoxic vs rest), the graph-informed model attains accuracy 0.864, balanced accuracy 0.856, macro-F1 0.850, MCC 0.706, AUROC 0.950, AUPRC 0.910. The six tabular baselines were evaluated on the \textit{identical} seed-42 folds (Table~\ref{tab:classification}), and we report paired Wilcoxon/t-tests of the GNN against each baseline on MCC and AUROC (\texttt{model\_comparison\_stats.tsv}).

The comparison yields a clear result: the graph-informed model is \textbf{statistically on par with the strongest gradient-boosting baselines} -- neither LightGBM (MCC 0.733) nor XGBoost (0.727) differs significantly from the GNN (0.706) on paired tests ($\Delta$MCC $-$0.028, t-test p=0.28; and $-$0.022, p=0.31, respectively) -- and it \textbf{significantly outperforms} the linear, kernel, and shallow-network baselines (vs ElasticNet $\Delta$MCC +0.035, p=0.017; vs SVM +0.217, p=0.008; vs MLP +0.330, p=0.037; Figure~\ref{fig:model}). We therefore do not claim state-of-the-art accuracy; the graph model matches top tree ensembles on this binary task while providing a mechanism-structured gene embedding (used for the axis analyses and target prioritization) that the tabular baselines do not. This "comparable accuracy, added interpretability" position is the straightforward reading of the numbers.

\begin{table*}[htbp]
\caption{NK-state classification (TCGA-STAD, 5-fold CV; mean over folds).}\label{tab:classification}
\scriptsize
\resizebox{\textwidth}{!}{\begin{tabular}{@{}p{1.0cm}p{1.92cm}p{1.92cm}p{1.92cm}p{1.92cm}p{1.92cm}p{1.92cm}@{}}
\toprule
Method & Accuracy & Balanced Acc. & Macro F1 & MCC & AUROC & AUPRC \\
\midrule
LightGBM & 0.880 & 0.863 & 0.865 & \textbf{0.733} & \textbf{0.960} & 0.933 \\
XGBoost & 0.878 & 0.859 & 0.862 & 0.727 & 0.959 & 0.928 \\
\textbf{GC-NKGraph-Atlas (GNN)} & 0.864 & 0.856 & 0.850 & 0.706 & 0.950 & 0.910 \\
RandomForest & 0.860 & 0.825 & 0.836 & 0.681 & 0.941 & 0.901 \\
ElasticNet & 0.856 & 0.822 & 0.832 & 0.671 & 0.927 & 0.884 \\
SVM (RBF) & 0.773 & 0.672 & 0.683 & 0.489 & 0.900 & 0.839 \\
MLP (2-layer) & 0.747 & 0.669 & 0.641 & 0.376 & 0.807 & 0.730 \\
\bottomrule
\end{tabular}}
\end{table*}

\textit{GNN not significantly different from LightGBM/XGBoost (paired p>0.27); significantly above ElasticNet/SVM/MLP (paired p<0.05).}

\begin{figure*}[htbp]
\centering
\includegraphics[width=\textwidth]{figures/fig4_model_comparison.pdf}
\caption{NK-state classification performance (TCGA-STAD, 5-fold stratified CV): MCC and AUROC for GC-NKGraph-Atlas (GNN) versus six tabular baselines on identical seed-42 folds; error bars show cross-fold standard deviation. The GNN is statistically on par with LightGBM/XGBoost and significantly above ElasticNet/SVM/MLP. \textit{Alt text: Bar charts of MCC and AUROC comparing the graph model with six tabular baselines on identical cross-validation folds.}}\label{fig:model}
\end{figure*}

\subsection{Candidate target prioritization (de-circularized)}

The naive composite score is circular: weighting axis membership at 40\% while scoring |tumor-specificity| promotes NK-effector markers (NKG7, PRF1, GZMB), which are \textit{depleted} in malignant cells and are the axis \textit{readout}, not targets. We therefore separate two tables (\texttt{src/interpretation/split\_target\_lists.py}):

\begin{enumerate}
\item \textbf{Tumor-intrinsic candidates} (Table~\ref{tab:targets}; n=37): genes required to be \textit{up in malignant cells} (tumor\_specificity\_log2>0), re-scored on signed tumor-specificity, mechanistic centrality (tumor-serine program / axis-druggable enzyme), druggability, and NK association. The top of the list is dominated by the mechanistically-privileged, druggable head of the axis -- \textbf{PHGDH} (serine synthesis; Phase 1/2 inhibitor), \textbf{SGMS2/SMPD3/SMPD1} (SM synthesis/catabolism; the exact enzymes the anchor mechanism implicates), \textbf{PSAT1/PSPH} (serine pathway) -- alongside gastric-relevant surface targets (ERBB2, FGFR2, MET) and stress ligands (MICA).
\item \textbf{Axis-confirmation panel} (n=36): the NK-side markers (GZMB, GNLY, PRF1, NKG7, protrusion genes), reported explicitly as \textit{readout} validating the axis score, not as targets.
\end{enumerate}

The de-circularization is conservative but not complete: an audit of the n=37 tumor-intrinsic pool reveals that 17 genes (46\%) are annotated to NK-side mechanism-card modules (protrusion machinery 7, de novo sphingolipid 5, SM catabolism 4, SM synthesis 1), including the NK protrusion GTPase RAC1 (rank 10) and the actin nucleation factor WASL (rank 24). These genes appear in the pool because their expression is detectably above zero in malignant cells (satisfying \texttt{tumor\_specificity\_log2}>0), but their mechanism-card module membership means the ranking retains a structural NK-side bias. The top of the list -- PHGDH, SGMS2, PSAT1, PSPH, SMPD3/1 -- is clean (tumor-serine-capacity or metabolic-suppression categories) and these remain the primary candidates for experimental follow-up.

\begin{table*}[htbp]
\caption{Top putative tumor-intrinsic candidate targets in gastric cancer (excerpt). The modest log2 fold-change values reflect the conservative positive-selection threshold; statistical confidence intervals are reported in the full supplementary table.}\label{tab:targets}
\scriptsize
\resizebox{\textwidth}{!}{\begin{tabular}{@{}p{1.0cm}p{2.30cm}p{2.30cm}p{2.30cm}p{2.30cm}p{2.30cm}@{}}
\toprule
Rank & Gene & Category & tumor\_spec (log2FC) & Druggability & Recommended assay \\
\midrule
1 & PHGDH & serine metabolism & +0.059 & Phase 1/2 (PHGDH inhibitor) & NK--tumor co-culture + serine-pathway inhibition \\
2 & SGMS2 & SM synthesis & +0.038 & Preclinical & SM-synthase modulation + NK membrane-SM readout \\
3 & PSAT1 & serine metabolism & +0.016 & Preclinical & NK--tumor co-culture + serine deprivation \\
4 & PSPH & serine metabolism & +0.010 & Preclinical & NK--tumor co-culture + serine deprivation \\
5 & SMPD3 & SM catabolism & +0.032 & Preclinical (Zheng 2023) & SMase-inhibitor rescue $\pm$ Tim3 blockade \\
8 & SMPD1 & SM catabolism & +0.001 & Preclinical (Zheng 2023) & SMase-inhibitor rescue $\pm$ Tim3 blockade \\
… & ERBB2, FGFR2, MET, MICA & gastric / stress-ligand & + & mixed & IHC/qPCR in GC tissue + NK co-culture \\
\bottomrule
\end{tabular}}
\end{table*}

\begin{figure*}[htbp]
\centering
\includegraphics[width=\textwidth]{figures/fig3_targets.pdf}
\caption{(A) tumor-specificity vs NK-association scatter, putative tumor-intrinsic pool highlighted; (B) axis positions of the top candidates on the serine$\to$SM$\to$protrusion$\to$cytotoxicity chain; (C) druggability distribution. \textit{Alt text: Multi-panel target-prioritization figure showing tumor-specificity versus NK-association, candidate positions on the serine-SM-protrusion-cytotoxicity chain, and druggability categories.}}\label{fig:targets}
\end{figure*}

\subsection{Graph-edge ablation}

To test whether the mechanism-specific \texttt{metabolic\_crosstalk} edge contributes structure beyond generic network priors, we compared three graph variants by their spectral gene embedding quality (Methods \S{}2.6, Stage~1):

\begin{table*}[htbp]
\centering
\caption{Graph-edge ablation: spectral embedding quality under edge removal. Modularity = intra-module minus inter-module cosine similarity.}\label{tab:ablation}
\footnotesize
\begin{tabularx}{\textwidth}{@{}lXXXXX@{}}
\toprule
Variant & Edges & Modularity & H2 (SM$\leftrightarrow$protrusion) & H1 (serine$\leftrightarrow$SM) \\
\midrule
FULL (all edges) & 1,129 & 0.380 & 0.591 & 0.239 \\
$-$MC (without \texttt{metabolic\_crosstalk}) & 829 & 0.159 & 0.487 & 0.000 \\
$-$SST (without all SST edges) & 9 & 0.017 & 0.000 & 0.000 \\
\bottomrule
\end{tabularx}
\end{table*}

Removing the \texttt{metabolic\_crosstalk} edge (300 edges, weight 0.5) reduces modularity by 2.4-fold ($\Delta$=$-$0.222) and eliminates the tumor-serine$\leftrightarrow$NK-SM coupling entirely (H1: 0.239$\to$0.000). When both SST-specific edge types are removed, only 9 coexpression edges remain, producing an essentially random embedding (modularity 0.017). The \texttt{metabolic\_crosstalk} edge therefore measurably shapes the embedding -- the graph design is not cosmetic, and the mechanism-grounded edge has a confirmable structural effect beyond what generic PPI/ligand--receptor edges provide. This ablation adds empirical support to the paper's core architectural claim: the heterogeneous graph's value lies in its mechanism-structured embedding, not in raw classification accuracy (where the GNN is on par with LightGBM/XGBoost; \S{}3.4). The value of typed ligand--receptor and mechanism-grounded edges over generic co-expression is consistent with single-cell studies in which ligand--receptor interaction analysis, rather than expression alone, explained how one cell subset regulates another [46].

\section{Discussion}

\subsection{Summary of findings}

Our central result is a \textbf{map of the mechanism's transcriptional reach}, not a blanket claim of recovery. Three findings anchor it:

\begin{enumerate}
\item \textbf{The effector layer of the axis is transcriptionally recoverable and generalizes.} The protrusion-machinery$\to$cytotoxicity coupling replicates from independent public liver transcriptomes (r=0.55) and single NK cells (r=0.32), and the same coupling holds in gastric cancer -- extending the recoverable layer to a digestive-tract cancer not on the mechanism's published list.
\item \textbf{Cell-type resolution is not optional for the metabolic layer.} The SM-balance$\to$protrusion coupling is invisible in bulk and appears only in isolated NK cells. This is a concrete, quantitative demonstration of why the framework's single-cell attribution -- rather than bulk deconvolution -- is necessary. It is consistent with the \texttt{metabolic\_crosstalk} edge's design premise, but given the negligible magnitude of the recovered coupling (r\textsuperscript{2}$<$0.001) it does not by itself validate that premise.
\item \textbf{Transcription does not substitute for the physical topology phenotype.} Machinery-gene transcription runs \textit{opposite} to the physical protrusion collapse in intratumoral NK (higher transcript, lower function). This is the natural boundary of any transcriptome-based reconstruction of a membrane-lipid mechanism, and it precisely reproduces why the anchor lab required single-cell mass spectrometry and super-resolution imaging in the first place.
\end{enumerate}

Together these delineate which layers a transcriptome can and cannot reach -- a result of direct use to any lab extending a physical immune-evasion mechanism to cohort scale. The prioritized putative tumor-intrinsic target list, led by the druggable serine/sphingomyelin enzymes at the mechanistic head of the axis, converts this map into experimentally testable follow-up. Additionally, a per-sample readout combining NK SM-catabolism score and HAVCR2 (Tim3) expression can stratify samples by the logic of the SM-restoration + Tim3-blockade combination proposed in the anchor paper -- presented as an in-silico hypothesis for experimental testing, not a validated clinical predictor. Consistent with this checkpoint-combination logic, inhibitory receptors are co-regulated on tumor NK cells: an independent HCC study reports that a distinct NK checkpoint, CLEC12B, correlates in TCGA-HCC with HAVCR2/TIM-3, TIGIT, PDCD1 and LAG-3 [48], supporting the biological coherence of a combined SM-catabolism plus Tim3-axis readout.

\subsection{The mechanism-card approach}

A key contribution of this work is the mechanism-card abstraction itself -- a machine-readable YAML formalism that separates the computational operationalization of a mechanism from the pipeline that executes it. A card declares, for one published mechanism, its molecular chain, the gene modules and expected directions for each step, the cell type each module is attributed to, the physical ground-truth measurements that are \textit{out of scope} for transcriptome analysis, the graph node/edge types it introduces, and the pre-registered validation hypotheses with an explicit recovery definition. Because the pipeline consumes the card rather than hard-coding the biology, applying the framework to a new mechanism -- a different metabolic immune checkpoint (e.g. adenosine-mediated NK suppression), a different effector cell type, or a stress-ligand-shedding axis -- requires only authoring a new card, not rewriting the pipeline core. The card registry (\texttt{configs/mechanism\_cards/registry.yaml}) is designed to accumulate such cards over time.

Beyond reuse, the card format enforces scientific discipline that this study shows to be load-bearing rather than decorative. First, it requires \textbf{explicit claim boundaries} in the card body, which is what let us state without ambiguity that machinery transcription is not a proxy for the physical topology phenotype -- a boundary our own results confirmed empirically (\S{}3.2). Second, it \textbf{gates physical ground-truth targets} (SEM protrusion density, single-cell SM mass spectrometry) as a separate, non-transcriptional layer, preventing their accidental conflation with the computable proxy. Third, it requires \textbf{pre-registered hypotheses with a defined recovery criterion}, so that a partial or negative outcome (as in Arm A) is reported as a structured result rather than quietly reframed. In our experience, this discipline is precisely what makes a partial-recovery finding credible to a mechanistic wet lab, and it generalizes to any attempt to operationalize a physical mechanism from an indirect molecular readout.

\subsection{Limitations}

\begin{enumerate}
\item \textbf{Transcriptional proxy $\neq$ physical topology.} Gene expression captures the molecular machinery and capacity for the serine--SM--topology axis, not the physical membrane phenotype itself. The disciplined qualifiers throughout are load-bearing: every claim is bounded by "transcriptional program permissive-of / associated-with."
\item \textbf{Serine/SM crosstalk is a metabolite-level effect.} Transcription captures enzyme abundance, not flux. The actual serine$\to$SM crosstalk requires metabolomics or the anchor lab's single-cell mass spectrometry for direct measurement.
\item \textbf{Data availability constraints.} The anchor paper did not deposit transcriptomic data, requiring the use of independent public cohorts. The liver positive control is therefore not a direct replication but an independent validation.
\item \textbf{No experimental validation.} All targets are computationally prioritized; none have been tested in wet-lab assays. The recommended assays are offered as a bridge to experimental follow-up.
\item \textbf{NK subtype resolution.} scRNA-based NK annotation depends on the quality of the reference atlas. Populations that are rare or absent in the reference may be misclassified.
\item \textbf{The graph model does not outperform top tabular baselines on accuracy.} On the binary NK-state task the GNN is statistically indistinguishable from LightGBM/XGBoost (\S{}3.4). Its value here is the mechanism-structured gene embedding it produces for the axis and target analyses, not a raw predictive-accuracy gain; a task where relational structure is more decisive (e.g. multi-class NK-state or cross-cohort transfer) would test the architecture more stringently.
\item \textbf{Residual NK bias in the tumor-intrinsic candidate pool.} The \texttt{tumor\_specificity\_log2}>0 gate is permissive: 17 of 37 candidates are annotated to NK-side mechanism-card modules, including RAC1 and WASL. The pool should be interpreted as ``genes with a non-zero malignant-cell transcript signal that mechanistically intersect the SST axis,'' not as a clean set of tumor-exclusive targets. A stricter filter (e.g.\ tumor\_specificity\_log2>0.5 or a module-level penalty) would reduce NK-side contamination at the cost of losing borderline tumor-intrinsic candidates.
\item \textbf{Candidate atlas omits intracellular/TF-level regulators.} The prioritization scores surface and metabolic-enzyme genes; transcription-factor and intracellular negative regulators of NK function (e.g. the CREM/PKA--CREB axis [47]) fall outside the current candidate space and are a natural extension of the mechanism-card modules.
\item \textbf{No clinical-outcome anchor.} NK states are linked to scRNA-defined labels, not to patient survival or therapy response. Single-cell atlases that tie a functional state to durable clinical outcome [46] indicate a clear next step: anchoring the NK-state readout to outcome in a cohort with follow-up.
\end{enumerate}

\subsection{Future directions}

\begin{itemize}
\item \textbf{Additional cancer types.} Once the gastric extension is validated, the framework can test the axis in other digestive-tract cancers (colorectal, pancreatic, esophageal) using the same mechanism card.
\item \textbf{Additional mechanisms.} The mechanism-card registry (\texttt{configs/mechanism\_cards/registry.yaml}) is designed to hold multiple cards. Cards for adenosine-mediated NK suppression, TGF$\beta$-driven NK exclusion, stress-ligand shedding (MICA/B-ADAM17), and other NK checkpoint axes (TIGIT/CD96 [41,42]; the CLEC12B--lipoprotein-lipase axis recently shown to restrain tumor NK cells and to synergize with PD-1 blockade [48]) are natural next additions.
\item \textbf{Physical topology integration.} When membrane protrusion / microvilli imaging data become available (even for a subset of samples), Layer 14R-B of the SST-axis module can be activated to provide direct phenotype-transcriptome correlation.
\item \textbf{Prospective validation cohort.} A dedicated gastric cancer cohort with paired bulk RNA-seq, scRNA-seq, and functional NK assays would provide the strongest validation of the prioritized targets.
\end{itemize}

\section{Conclusion}

GC-NKGraph-Atlas demonstrates \textit{how far} a specific, published immune-evasion mechanism -- the serine--sphingomyelin--membrane-topology axis of NK dysfunction -- can be surveyed at cohort scale from public transcriptomes. Using a single-cell-informed heterogeneous graph framework and a two-arm design, we show that the effector layer of the axis recovers robustly and generalizes from liver to gastric cancer, that the upstream metabolic coupling recovers only when cell-type is resolved, and that the physical topology phenotype lies beyond the reach of machinery transcription. This scoping map, not an overclaim of full reconstruction, is the contribution, and it is exactly the kind of principled result that a mechanistic wet lab can build on. The mechanism-card abstraction is designed to make the framework a reusable engine rather than a one-off pipeline -- so that each published mechanism can be turned into a scalable target-discovery run whose transcriptional reach is measured, not assumed. We demonstrate this abstraction on a single card here; reuse across multiple mechanisms remains a design aspiration to be established as further cards are authored.

\section*{Data Availability}

All code and configuration files are available at \url{https://github.com/nblvguohao/GC-NKGraph-Atlas}. The mechanism-card template and the Zheng 2023 NK SM-topology card are provided under \texttt{configs/mechanism\_cards/}. Synthetic test data can be generated via \texttt{python src/common/synthetic\_data.py}.

\textbf{Real data availability:}

\begin{itemize}
\item TCGA-STAD and TCGA-LIHC: available from the Genomic Data Commons (\url{https://portal.gdc.cancer.gov/})
\item GSE62254, GSE84437: available from the Gene Expression Omnibus (\url{https://www.ncbi.nlm.nih.gov/geo/})
\item No novel sequencing data were generated for this study.
\end{itemize}

\textbf{Pre-registration:} The hypotheses H1--H5 and the original full-recovery criterion for the positive-control arm were registered in \texttt{configs/sst\_axis\_config.yaml} before execution. Because the full criterion was not met, the manuscript reports the result as a partial-recovery scoping map and explicitly separates pre-registered outcomes from post-hoc interpretation.

\section*{Author Contributions}

\textbf{Guohao Lyu:} Conceptualization, Methodology, Software, Formal analysis, Investigation, Data curation, Visualization, Writing -- original draft. \textbf{Yingchun Xia:} Methodology, Software, Validation, Writing -- review \& editing. \textbf{Huichao Liu:} Software, Data curation, Validation. \textbf{Xiaolei Zhu:} Formal analysis, Methodology, Writing -- review \& editing. \textbf{Shuai Yang:} Investigation, Validation, Visualization. \textbf{Ailian Zhou:} Conceptualization, Supervision, Funding acquisition, Writing -- review \& editing. \textbf{Lichuan Gu:} Conceptualization, Supervision, Project administration, Funding acquisition, Writing -- review \& editing.

All authors read and approved the final manuscript. (CRediT taxonomy; adjust individual roles to reflect actual contributions before submission.)

\section*{Ethics Approval and Consent to Participate}

Not applicable. This study used only publicly available, de-identified datasets (TCGA-LIHC, TCGA-STAD, GSE62254, GSE84437, GSE246662) obtained from the Genomic Data Commons and the Gene Expression Omnibus. No new human or animal subjects were involved, and no identifiable personal data were generated or analyzed. Use of these public data complies with the respective data-access policies.

\section*{Competing Interests}

The authors declare that they have no competing interests.

\section*{Funding}

This work was supported by grants from the National Natural Science Foundation of China (32472007, 62301006, 62301008), the Natural Science Foundation of Anhui Province (2308085MF217, 2308085QF202), and the Anhui Province Key Laboratory of Intelligent Agricultural Technology and Equipment.


\begin{thebibliography}{99}
\bibitem[1]{ref1} Vivier E, Tomasello E, Baratin M, et al. Functions of natural killer cells. \textit{Nat Immunol} 2008;9:503--10.
\bibitem[2]{ref2} Chiossone L, Dumas PY, Vienne M, Vivier E. Natural killer cells and other innate lymphoid cells in cancer. \textit{Nat Rev Immunol} 2018;18:671--88.
\bibitem[3]{ref3} O'Brien KL, Finlay DK. Immunometabolism and natural killer cell responses. \textit{Nat Rev Immunol} 2019;19:282--90.
\bibitem[4]{ref4} Terrén I, Orrantia A, Vitallé J, et al. NK cell metabolism and tumor microenvironment. \textit{Front Immunol} 2019;10:2278.
\bibitem[5]{ref5} André P, Denis C, Soulas C, et al. Anti-NKG2A mAb is a checkpoint inhibitor that promotes anti-tumor immunity by unleashing both T and NK cells. \textit{Cell} 2018;175:1731--43.e13.
\bibitem[6]{ref6} Gao J, Zheng X, Liu Y, et al. TGF-$\beta$ impairs NK cell migration and cytotoxicity. \textit{J Immunother Cancer} 2023;11:e005785.
\bibitem[7]{ref7} Melaiu O, Lucarini V, Cifaldi L, Fruci D. Influence of the tumor microenvironment on NK cell function in solid tumors. \textit{Front Immunol} 2020;10:3038.
\bibitem[8]{ref8} Zheng X, Hou Z, Qian Y, et al. Tumors evade immune cytotoxicity by altering the surface topology of NK cells. \textit{Nat Immunol} 2023;24:748--59. doi:10.1038/s41590-023-01462-9.
\bibitem[9]{ref9} Zheng X, Tian Z, Wei H. Sphingomyelin is a prospective metabolic immune checkpoint for natural killer cells. \textit{Clin Transl Med} 2023;13:e1395.
\bibitem[10]{ref10} Hu Z, Dong Y, Wang K, Sun Y. Heterogeneous graph transformer. \textit{Proc Web Conf} 2020:2704--10.
\bibitem[11]{ref11} Chen T, Guestrin C. XGBoost: a scalable tree boosting system. \textit{Proc KDD} 2016:785--94.
\bibitem[12]{ref12} Ke G, Meng Q, Finley T, et al. LightGBM: a highly efficient gradient boosting decision tree. \textit{Adv NeurIPS} 2017:3146--54.
\bibitem[13]{ref13} Wolf FA, Angerer P, Theis FJ. SCANPY: large-scale single-cell gene expression data analysis. \textit{Genome Biol} 2018;19:15.
\bibitem[14]{ref14} Gayoso A, Lopez R, Xing G, et al. A Python library for probabilistic analysis of single-cell omics data. \textit{Nat Biotechnol} 2022;40:163--66.
\bibitem[15]{ref15} Fey M, Lenssen JE. Fast graph representation learning with PyTorch Geometric. \textit{ICLR Workshop} 2019.
\bibitem[16]{ref16} Efremova M, Vento-Tormo M, Teichmann SA, Vento-Tormo R. CellPhoneDB: inferring cell--cell communication from combined expression of multi-subunit ligand--receptor complexes. \textit{Nat Protoc} 2020;15:1484--506.
\bibitem[17]{ref17} Newman AM, Steen CB, Liu CL, et al. Determining cell type abundance and expression from bulk tissues with digital cytometry. \textit{Nat Biotechnol} 2019;37:773--82.
\bibitem[18]{ref18} Finotello F, Mayer C, Plattner C, et al. Molecular and pharmacological modulators of the tumor immune contexture revealed by deconvolution of RNA-seq data. \textit{Genome Med} 2019;11:34.
\bibitem[19]{ref19} Sun D, Guan X, Moran AE, et al. Identifying phenotype-associated subpopulations by integrating bulk and single-cell sequencing data. \textit{Nat Biotechnol} 2022;40:527--38.
\bibitem[20]{ref20} Jin S, Guerrero-Juarez CF, Zhang L, et al. Inference and analysis of cell--cell communication using CellChat. \textit{Nat Commun} 2021;12:1088.
\bibitem[21]{ref21} Browaeys R, Saelens W, Saeys Y. NicheNet: modeling intercellular communication by linking ligands to target genes. \textit{Nat Methods} 2020;17:159--62.
\bibitem[22]{ref22} Alghamdi N, Chang W, Dang P, et al. A graph neural network model to estimate cell-wise metabolic flux using single-cell RNA-seq data. \textit{Genome Res} 2021;31:1867--84.
\bibitem[23]{ref23} Gulati GS, Sikandar SS, Wesche DJ, et al. Single-cell transcriptional diversity is a hallmark of developmental potential. \textit{Science} 2020;367:405--11.
\bibitem[24]{ref24} Wang T, Shao W, Huang Z, et al. MOGONET integrates multi-omics data using graph convolutional networks allowing patient classification and biomarker identification. \textit{Nat Commun} 2021;12:3445.
\bibitem[25]{ref25} Huntington ND, Cursons J, Rautela J. The cancer--natural killer cell immunity cycle. \textit{Nat Rev Cancer} 2020;20:437--54.
\bibitem[26]{ref26} Myers JA, Miller JS. Exploring the NK cell platform for cancer immunotherapy. \textit{Nat Rev Clin Oncol} 2021;18:85--100.
\bibitem[27]{ref27} Shimasaki N, Jain A, Campana D. NK cells for cancer immunotherapy. \textit{Nat Rev Drug Discov} 2020;19:200--18.
\bibitem[28]{ref28} Possemato R, Marks KM, Shaul YD, et al. Functional genomics reveal that the serine synthesis pathway is essential in breast cancer. \textit{Nature} 2011;476:346--50.
\bibitem[29]{ref29} Locasale JW. Serine, glycine and one-carbon units: cancer metabolism in full circle. \textit{Nat Rev Cancer} 2013;13:572--83.
\bibitem[30]{ref30} Ogretmen B. Sphingolipid metabolism in cancer -- signalling and drug resistance. \textit{Nat Rev Cancer} 2018;18:33--50.
\bibitem[31]{ref31} Hannun YA, Obeid LM. Sphingolipids and their metabolism in physiology and disease. \textit{Nat Rev Mol Cell Biol} 2018;19:175--91.
\bibitem[32]{ref32} Hao Y, Hao S, Andersen-Nissen E, et al. Integrated analysis of multimodal single-cell data. \textit{Cell} 2021;184:3573--87.e29.
\bibitem[33]{ref33} Stuart T, Butler A, Hoffman P, et al. Comprehensive integration of single-cell data. \textit{Cell} 2019;177:1888--902.
\bibitem[34]{ref34} Luecken MD, Theis FJ. Current best practices in single-cell RNA-seq analysis: a tutorial. \textit{Mol Syst Biol} 2019;15:e8746.
\bibitem[35]{ref35} Lotfollahi M, Naghipourfar M, Luecken MD, et al. Mapping single-cell data to reference atlases by transfer learning. \textit{Nat Biotechnol} 2022;40:121--30.
\bibitem[36]{ref36} Kipf TN, Welling M. Semi-supervised classification with graph convolutional networks. \textit{Proc ICLR} 2017.
\bibitem[37]{ref37} Veličković P, Cucurull G, Casanova A, et al. Graph attention networks. \textit{Proc ICLR} 2018.
\bibitem[38]{ref38} Schlichtkrull M, Kipf TN, Bloem P, et al. Modeling relational data with graph convolutional networks. \textit{Proc ESWC} 2018:593--607.
\bibitem[39]{ref39} Cancer Genome Atlas Research Network. Comprehensive molecular characterization of gastric adenocarcinoma. \textit{Nature} 2014;513:202--9.
\bibitem[40]{ref40} Cristescu R, Lee J, Nebozhyn M, et al. Molecular analysis of gastric cancer identifies subtypes associated with distinct clinical outcomes. \textit{Nat Med} 2015;21:449--56.
\bibitem[41]{ref41} Barry KC, Hsu J, Broz ML, et al. A natural killer--dendritic cell axis defines checkpoint therapy--responsive tumor microenvironments. \textit{Nat Med} 2018;24:1178--91.
\bibitem[42]{ref42} Zhang Q, Bi J, Zheng X, et al. Blockade of the checkpoint receptor TIGIT prevents NK cell exhaustion and elicits potent anti-tumor immunity. \textit{Nat Immunol} 2018;19:723--32.
\bibitem[43]{ref43} Subramanian A, Tamayo P, Mootha VK, et al. Gene set enrichment analysis: a knowledge-based approach for interpreting genome-wide expression profiles. \textit{Proc Natl Acad Sci USA} 2005;102:15545--50.
\bibitem[44]{ref44} Liberzon A, Birger C, Thorvaldsdóttir H, et al. The Molecular Signatures Database (MSigDB) hallmark gene set collection. \textit{Cell Syst} 2015;1:417--25.
\bibitem[45]{ref45} Laskowski TJ, Biederstädt A, Rezvani K. Natural killer cells in antitumour adoptive cell immunotherapy. \textit{Nat Rev Cancer} 2022;22:557--75.
\bibitem[46]{ref46} Bai Z, Feng B, McClory SE, et al. Single-cell CAR T atlas reveals type 2 function in 8-year leukaemia remission. \textit{Nature} 2024;634:702--11. doi:10.1038/s41586-024-07762-w.
\bibitem[47]{ref47} Rafei H, Basar R, Acharya S, et al. CREM is a regulatory checkpoint of CAR and IL-15 signalling in NK cells. \textit{Nature} 2025;643:1076--86. doi:10.1038/s41586-025-09087-8.
\bibitem[48]{ref48} Sun P, Xu X, Hu B, et al. Targeting NK cell CLEC12B enhances cancer immunotherapy. \textit{Nat Immunol} 2026;27:985--99. doi:10.1038/s41590-026-02471-0.
\end{thebibliography}

\end{document}
